\chapter{Introduction}
\label{chap:introduction}

The introduction should awaken the reader's interest in the contents of
this report, point out the treated problems, and describe the concepts that
were developed for their solution. For this purpose, this chapter first
explains the motivation of the project in Section~1.1, then defines and
delimits the problem in Section~1.2, fixes the objective of the work in
Section~1.3, and finally presents the chosen approach in Section~1.4.

\section{Motivation}
\label{sec:motivation}

In classical machine learning, it is assumed that the complete training
data set is available from the beginning and that all samples are
independent and identically distributed. Under this assumption, modern deep
neural networks achieve impressive results in image classification, speech
recognition and many other areas. In realistic applications, however, this
assumption is often not fulfilled. Data arrives step by step: a medical
diagnosis system sees new diseases, a recommendation system sees new user
groups, and a self driving car sees new traffic situations, weather
conditions or regions. In all these cases, the model has to be updated
continuously, ideally without storing and reprocessing all historical data
again.

A simple update strategy, namely to continue the training of the network on
the new data, fails in a drastic and well documented way. The network
adapts its weights to the new task and thereby overwrites exactly those
weight configurations that stored the solution of the earlier tasks. After
only a few update steps, the performance on the old data breaks down,
although the architecture and the capacity of the network have not changed.
This phenomenon is known as catastrophic forgetting or catastrophic
interference and has been studied in neural networks since the late 1980s
\cite{mccloskey1989,french1999}. It is still seen as one of the central
obstacles on the way to adaptive, lifelong learning systems
\cite{delange2021,wang2024survey}.

The research field of continual learning, also called lifelong or
incremental learning, develops methods against this problem. Roughly
speaking, three families of approaches exist. Regularisation based methods
such as Elastic Weight Consolidation \cite{kirkpatrick2017} punish changes
to weights that were important for earlier tasks. Replay based methods
store a small memory of old data and mix it into the training of new tasks
\cite{rebuffi2017,lopezpaz2017}. Architecture based methods change the
structure of the network itself so that new knowledge can be stored in
previously unused components \cite{rusu2016}. The present project belongs
to the third family.

The starting point of this work is the Mixture of Experts architecture.
Instead of one single large network, such a model consists of several
smaller sub networks, the so called experts, and a learned gating network,
also called router, that decides for every input which experts should
process it \cite{jacobs1991,shazeer2017}. The central hypothesis of this
project is that this structure is naturally suited for continual learning:
if the router learns to send the data of a new task to a different subset
of experts than the data of earlier tasks, then the gradient updates for
the new task mainly change those experts, and the knowledge stored in the
remaining experts stays mostly untouched. In other words, the model can
continue learning without overwriting, because the capacity of the model is
partitioned instead of being shared completely.

For me as a student, this topic is attractive because it connects two
important and modern topics of machine learning in a very practical way. On
the one hand, continual learning is a requirement for almost every machine
learning system that is used over a long time. On the other hand, Mixture
of Experts models are currently one of the most important building blocks
of very large language models \cite{fedus2022}, and it is an interesting
question whether their routing mechanism has useful properties beyond pure
scaling.

\section{Problem Statement and Delimitation}
\label{sec:problem}

A normal deep neural network works well when all training data is presented
at the same time. When the data arrives one task after another, the network
loses its old abilities quickly. The concrete problem of this project is
therefore to build and evaluate a system that fulfils three requirements at
the same time. First, plasticity: the system must be able to learn a new
task fast and with high accuracy, using only the data of the current task.
Second, stability: after a sequence of tasks, the system must still reach a
high accuracy on the earlier tasks, without being trained on their data
again. Third, efficiency: the additional computational cost and the
additional memory compared to a normal network must stay moderate.

These three requirements are in conflict with each other, which is known as
the stability plasticity dilemma \cite{mermillod2013}. A completely stable
model cannot learn anything new, and a completely plastic model forgets
everything old. The problem of this project is to find out whether a
Mixture of Experts architecture can offer a better compromise than a
monolithic network.

To keep the problem tractable, the project is delimited in the following
way. The investigation is restricted to two established benchmarks that are
built from the MNIST data set \cite{lecun1998}: SplitMNIST for the task
incremental scenario and PermutedMNIST for the domain incremental scenario
\cite{vandeven2019}. The harder class incremental scenario, in which the
model also has to infer the task identity of a test sample, is not part of
this work. The investigated networks are feed forward networks of moderate
size; convolutional architectures and larger image data sets are excluded.
Replay methods, which store old data, are deliberately not used, because
the goal is to find out how far a purely architectural mechanism can go
without any stored data. Extensions in these directions are discussed as
outlook in Chapter~8.

\section{Objective of the Work}
\label{sec:objective}

The objective of this work is to build and test a Mixture of Experts model
for continual learning and to find out whether it forgets less than a
normal neural network. Following the project proposal, this objective is
split into five concrete goals:

\begin{enumerate}[label=\textbf{G\arabic*.}, leftmargin=3em]
  \item Study why normal neural networks forget old data so quickly, and
        summarise the theoretical background of continual learning and
        Mixture of Experts models.
  \item Design and implement a Mixture of Experts model with a trainable
        router that can be trained step by step on a sequence of tasks.
  \item Train this model on two different benchmarks that present the data
        one task after another.
  \item Compare the model with simple neural networks and with an
        established regularisation method to show how much it forgets in
        comparison.
  \item Examine how the router works and how its behaviour contributes to
        keeping previously learned knowledge.
\end{enumerate}

The work is considered successful if all five goals are reached with
reproducible experiments and if the comparison is based on the standard
metrics of the field.

\section{Approach}
\label{sec:approach}

After problem statement and objective have fixed the starting point and the
end point of the work, this section presents the path that was taken. The
individual solution steps and their connection are shown in
Figure~1.1.

\begin{figure}[H]
\centering
\begin{tikzpicture}[
  pstep/.style={draw, rounded corners, minimum height=0.95cm, minimum width=2.65cm, align=center, fill=blue!6, font=\small},
  arr/.style={-{Stealth[length=2.5mm]}, thick}]
  \node[pstep] (s1) {1. Literature\\review};
  \node[pstep, right=0.55cm of s1] (s2) {2. Benchmarks\\SplitMNIST,\\PermutedMNIST};
  \node[pstep, right=0.55cm of s2] (s3) {3. Baselines\\naive, wide,\\EWC};
  \node[pstep, right=0.55cm of s3] (s4) {4. MoE model\\and router\\variants};
  \node[pstep, below=0.7cm of s3] (s5) {5. Experiments\\3 seeds, CL\\metrics};
  \node[pstep, right=0.55cm of s5] (s6) {6. Router\\analysis};
  \node[pstep, left=0.55cm of s5] (s7) {7. Evaluation\\and report};
  \draw[arr] (s1) -- (s2);
  \draw[arr] (s2) -- (s3);
  \draw[arr] (s3) -- (s4);
  \draw[arr] (s4.south) |- (s5.east);
  \draw[arr] (s5) -- (s6);
  \draw[arr] (s6) -- (s7);
\end{tikzpicture}
\caption{Overview of the approach of this project.}
\label{fig:approach}
\end{figure}

In the first step, the literature on catastrophic forgetting, continual
learning and Mixture of Experts models was studied, and the theoretical
foundation was worked out. This foundation is presented in Chapter~2. In
the second step, the two benchmarks were constructed and the problem was
analysed in detail; this analysis and the resulting requirements are
documented in Chapter~3. In the third and fourth step, the baseline models
and the Mixture of Experts model with its router were designed, which is
the content of Chapter~4. The concrete implementation in Python and
PyTorch is described in Chapter~5. Chapter~6 explains how the implemented
programs were brought into operation and how the experiment suite was
executed. Chapter~7 evaluates the results of all experiments, including a
detailed router analysis and an architecture variant that resulted from the
findings. Chapter~8 finally summarises the achieved results and gives an
outlook on possible extensions.
